\paragraph{同余块 Abel--Weil 展开、留数锁相与误差可审计重构链}\label{par:xi-discrete-abel-weil-congruence-residue-locking-audit}
在定理 \ref{thm:discrete-abel-weil-expansion} 的离散 Abel--Weil 分解基础上，
可将底数采样误差序列按同余类拆分为一组共享极点的子通道。
该拆分把“零点参数 \(\rho\)”转写为“公共极点位置 + 通道留数相位/模长”的代数可观测量，
并允许在观测误差下给出显式的非渐近传播上界。

\begin{theorem}[同余块 Abel--Weil 展开与公共极点留数锁相]\label{thm:xi-congruence-block-abel-weil-residue-locking}
固定 \(b>1\) 与整数 \(m\ge 2\)。
沿用
\[
E_b(n):=b^{-n/2}\bigl(\psi_0(b^n)-b^n\bigr),
\qquad
\mathcal A_b(r):=\sum_{n\ge 0}E_b(n)r^n.
\]
对每个 \(j\in\{0,1,\dots,m-1\}\) 定义同余块生成函数
\[
\mathcal A_{b,m}^{(j)}(u):=\sum_{n\ge 0}E_b(mn+j)\,u^n.
\]
记 \(\zeta\) 非平凡零点的重数为 \(\mathrm{mult}(\rho)\)。
则在 \(|u|\) 足够小（继而按亚纯延拓）时有
\begin{equation}\label{eq:xi-congruence-block-abel-weil}
\mathcal A_{b,m}^{(j)}(u)
=-\sum_{\rho}\frac{\mathrm{mult}(\rho)\,b^{(\rho-\frac12)j}}{\rho}\cdot
\frac{1}{1-u\,b^{m(\rho-\frac12)}}
+H_{b,m}^{(j)}(u),
\end{equation}
其中
\[
H_{b,m}^{(j)}(u):=\sum_{n\ge 0}b^{-(mn+j)/2}G(b^{mn+j})\,u^n
\]
至少在 \(|u|<b^{m/2}\) 内解析。
于是每个非平凡零点 \(\rho\) 诱导公共极点
\[
u_\rho:=b^{-m(\rho-\frac12)},
\]
且在第 \(j\) 通道的留数满足
\begin{equation}\label{eq:xi-congruence-block-residue-locking}
\operatorname{Res}_{u=u_\rho}\mathcal A_{b,m}^{(j)}(u)
=\frac{\mathrm{mult}(\rho)\,b^{(\rho-\frac12)j}\,u_\rho}{\rho}.
\end{equation}
\end{theorem}

\begin{proof}
由 \eqref{eq:Abel-Weil-decomposition} 的系数形式，
\[
E_b(n)= -\sum_\rho \frac{\mathrm{mult}(\rho)}{\rho}\,b^{(\rho-\frac12)n}
+b^{-n/2}G(b^n).
\]
代入 \(n=mn+j\) 并求和得
\[
\mathcal A_{b,m}^{(j)}(u)
=-\sum_\rho \frac{\mathrm{mult}(\rho)\,b^{(\rho-\frac12)j}}{\rho}
\sum_{n\ge 0}\bigl(u\,b^{m(\rho-\frac12)}\bigr)^n
+\sum_{n\ge 0}b^{-(mn+j)/2}G(b^{mn+j})u^n,
\]
即 \eqref{eq:xi-congruence-block-abel-weil}。
由 \(G(b^{mn+j})=O(1)\) 得余项系数 \(O(b^{-(mn+j)/2})\)，故解析半径至少为 \(b^{m/2}\)。
对主项的几何核
\(
(1-u\,b^{m(\rho-\frac12)})^{-1}
\)
在 \(u=u_\rho\) 处取留数即得 \eqref{eq:xi-congruence-block-residue-locking}。证毕。
\end{proof}

\begin{corollary}[留数比恢复零点实部与模长平坦判据]\label{cor:xi-congruence-block-residue-ratio-realpart}
在定理 \ref{thm:xi-congruence-block-abel-weil-residue-locking} 设定下，
若 \(\rho\) 为单零点，则定义
\[
r_\rho:=b^{-(\rho-\frac12)},\qquad u_\rho=r_\rho^m.
\]
对任意 \(j\in\{0,\dots,m-2\}\) 有
\begin{equation}\label{eq:xi-congruence-block-residue-ratio}
\frac{\operatorname{Res}_{u=u_\rho}\mathcal A_{b,m}^{(j)}}
{\operatorname{Res}_{u=u_\rho}\mathcal A_{b,m}^{(j+1)}}
=r_\rho.
\end{equation}
因此
\begin{equation}\label{eq:xi-congruence-block-realpart-from-residue-ratio}
\Re(\rho)-\frac12
=-\frac{1}{\log b}\log|r_\rho|
=-\frac{1}{\log b}
\log\left|
\frac{\operatorname{Res}\mathcal A_{b,m}^{(j)}}
{\operatorname{Res}\mathcal A_{b,m}^{(j+1)}}
\right|.
\end{equation}
并且
\[
\Re(\rho)=\frac12
\Longleftrightarrow
\left|\operatorname{Res}\mathcal A_{b,m}^{(j)}\right|
=\left|\operatorname{Res}\mathcal A_{b,m}^{(j+1)}\right|
\quad(0\le j\le m-2).
\]
\end{corollary}

\begin{proof}
由 \eqref{eq:xi-congruence-block-residue-locking}（\(\mathrm{mult}(\rho)=1\)）直接相除得到
\eqref{eq:xi-congruence-block-residue-ratio}。取模与对数即得
\eqref{eq:xi-congruence-block-realpart-from-residue-ratio}。证毕。
\end{proof}

\begin{proposition}[幂覆盖的解析截面由留数比内禀给出]\label{prop:xi-congruence-block-power-cover-section}
在推论 \ref{cor:xi-congruence-block-residue-ratio-realpart} 条件下，
对公共极点 \(u_\rho\) 定义
\[
s_m(u_\rho):=
\frac{\operatorname{Res}_{u=u_\rho}\mathcal A_{b,m}^{(j)}}
{\operatorname{Res}_{u=u_\rho}\mathcal A_{b,m}^{(j+1)}}
\qquad(0\le j\le m-2).
\]
则 \(s_m(u_\rho)\) 与 \(j\) 无关，且
\[
s_m(u_\rho)=r_\rho,\qquad s_m(u_\rho)^m=u_\rho.
\]
换言之，在零点诱导极点除子上，幂映射 \(\pi_m:z\mapsto z^m\) 存在由留数数据唯一确定的截面。
\end{proposition}

\begin{proof}
由 \eqref{eq:xi-congruence-block-residue-ratio} 得 \(s_m(u_\rho)=r_\rho\) 与 \(j\) 无关。
再由 \(u_\rho=r_\rho^m\) 即得 \(s_m(u_\rho)^m=u_\rho\)。证毕。
\end{proof}

\begin{theorem}[单基去混叠的可数例外集定理]\label{thm:xi-single-base-dealiasing-countable-exception}
\leanverified{paper\_xi\_single\_base\_dealiasing\_countable\_exception}
设 \(S\subset\CC\) 为可数集合。
对每个 \(b>1\) 定义
\[
\Phi_b:S\to\CC^\times,\qquad
\Phi_b(s):=b^{-(s-\frac12)}.
\]
则存在至多可数的例外集合 \(\mathcal B\subset(1,\infty)\)，使得对任意
\(
b\in(1,\infty)\setminus\mathcal B
\)，\(\Phi_b\) 在 \(S\) 上单射。
\end{theorem}

\begin{proof}
固定 \(s\neq s'\in S\)。
若 \(\Phi_b(s)=\Phi_b(s')\)，则
\[
b^{-(s-s')}=1
\Longleftrightarrow
(s-s')\log b\in 2\pi i\,\ZZ.
\]
若 \(\Re(s-s')\neq 0\)，上式无解；
若 \(\Re(s-s')=0\)，则对每个非零整数 \(k\) 至多给出一个解
\(
\log b=2\pi k/\Im(s-s')
\)。
故每对 \((s,s')\) 的碰撞底数集合至多可数。
对可数对集合 \(S\times S\setminus\{(s,s)\}\) 取并，得到例外集 \(\mathcal B\) 至多可数。证毕。
\end{proof}

\begin{theorem}[双基残差比编码的唯一重构]\label{thm:xi-two-base-residue-ratio-unique-recovery}
\leanverified{paper\_xi\_two\_base\_residue\_ratio\_unique\_recovery}
取 \(b_1,b_2>1\) 满足
\[
\frac{\log b_1}{\log b_2}\notin\QQ.
\]
设 \(\rho\) 为非平凡单零点，并在两底数下定义
\[
r_\rho^{(k)}
:=b_k^{-(\rho-\frac12)}
=\frac{\operatorname{Res}\mathcal A_{b_k,m}^{(j)}}
{\operatorname{Res}\mathcal A_{b_k,m}^{(j+1)}},
\qquad k=1,2.
\]
则有序对 \(\bigl(r_\rho^{(1)},r_\rho^{(2)}\bigr)\) 唯一决定 \(\rho\)。
\end{theorem}

\begin{proof}
若 \(\rho,\rho'\) 诱导同一编码对，则
\[
b_1^{-(\rho-\rho')}=1,\qquad b_2^{-(\rho-\rho')}=1.
\]
故存在整数 \(k_1,k_2\) 使
\[
\rho-\rho'=\frac{2\pi i k_1}{\log b_1}
=\frac{2\pi i k_2}{\log b_2}.
\]
若 \(\rho\neq\rho'\)，则
\(
\log b_1/\log b_2=k_1/k_2\in\QQ
\)，与假设矛盾，故 \(\rho=\rho'\)。证毕。
\end{proof}

\begin{proposition}[留数观测误差到零点实部误差的显式传播界]\label{prop:xi-residue-ratio-error-propagation-realpart}
固定 \(b>1,m\ge2\)。
设在公共极点 \(u_\rho\) 处，相邻通道真留数为 \(R_j,R_{j+1}\neq 0\)，观测值为
\(\widehat R_j,\widehat R_{j+1}\)，且
\[
\left|\frac{\widehat R_j}{R_j}-1\right|\le\varepsilon,\qquad
\left|\frac{\widehat R_{j+1}}{R_{j+1}}-1\right|\le\varepsilon,\qquad
0\le\varepsilon<1.
\]
定义
\[
r:=\frac{R_j}{R_{j+1}}=b^{-(\rho-\frac12)},
\qquad
\widehat r:=\frac{\widehat R_j}{\widehat R_{j+1}}.
\]
则
\begin{equation}\label{eq:xi-residue-ratio-relative-error-bound}
\left|\frac{\widehat r}{r}-1\right|
\le
\frac{2\varepsilon}{1-\varepsilon}.
\end{equation}
并且若
\(
\widehat\rho:=\frac12-\frac{\Log(\widehat r)}{\log b}
\)
（取一致对数分支），则
\begin{equation}\label{eq:xi-residue-ratio-realpart-error-bound}
\left|\Re(\widehat\rho)-\Re(\rho)\right|
\le
\frac{1}{\log b}
\log\!\left(1+\frac{2\varepsilon}{1-\varepsilon}\right).
\end{equation}
\end{proposition}

\begin{proof}
写
\(
\widehat R_j=R_j(1+\delta_j),\ 
\widehat R_{j+1}=R_{j+1}(1+\delta_{j+1})
\)
且 \(|\delta_j|,|\delta_{j+1}|\le\varepsilon\)。
则
\[
\frac{\widehat r}{r}
=\frac{1+\delta_j}{1+\delta_{j+1}},
\qquad
\left|\frac{\widehat r}{r}-1\right|
=\left|\frac{\delta_j-\delta_{j+1}}{1+\delta_{j+1}}\right|
\le\frac{|\delta_j|+|\delta_{j+1}|}{1-\varepsilon}
\le \frac{2\varepsilon}{1-\varepsilon},
\]
得 \eqref{eq:xi-residue-ratio-relative-error-bound}。
再由
\[
\Re(\widehat\rho)-\Re(\rho)
=-\frac{1}{\log b}\log\frac{|\widehat r|}{|r|},
\]
与不等式 \(|\log(1+z)|\le \log(1+|z|)\)（当 \(|z|<1\)）及
\eqref{eq:xi-residue-ratio-relative-error-bound}，得
\eqref{eq:xi-residue-ratio-realpart-error-bound}。证毕。
\end{proof}

\begin{theorem}[单极点轮廓积分读出与可审计误差界]\label{thm:xi-single-pole-contour-readout-error-audit}
设 \(F\) 在闭圆盘 \(\overline{D(0,R)}\) 邻域内亚纯，并在 \(D(0,R)\) 内恰有一个简单极点 \(u_0\)。
令 \(\Gamma_R:=\{u:\ |u|=R\}\)，定义
\[
J_0:=\oint_{\Gamma_R}F(u)\,\dd u,\qquad
J_1:=\oint_{\Gamma_R}uF(u)\,\dd u.
\]
则有精确读出
\[
\operatorname{Res}_{u=u_0}F=\frac{1}{2\pi i}J_0,\qquad
u_0=\frac{J_1}{J_0}.
\]
若 \(\widetilde F\) 是 \(\Gamma_R\) 上近似，且
\[
\sup_{|u|=R}|F(u)-\widetilde F(u)|\le \eta,
\]
定义 \(\widetilde J_k:=\oint_{\Gamma_R}u^k\widetilde F(u)\,\dd u\)（\(k=0,1\)），并假设非退化条件
\[
|\widetilde J_0|>2\pi R\eta,
\]
则 \(J_0\neq 0\)，并有
\begin{equation}\label{eq:xi-contour-residue-error-bound}
\left|\operatorname{Res}_{u=u_0}F-\frac{1}{2\pi i}\widetilde J_0\right|
\le R\eta,
\end{equation}
\begin{equation}\label{eq:xi-contour-pole-location-error-bound}
\left|u_0-\frac{\widetilde J_1}{\widetilde J_0}\right|
\le
\frac{R^2\eta}{\frac{|\widetilde J_0|}{2\pi}-R\eta}
+
\frac{R^3\eta}{\left(\frac{|\widetilde J_0|}{2\pi}-R\eta\right)\frac{|\widetilde J_0|}{2\pi}}
\cdot
\sup_{|u|=R}|\widetilde F(u)|.
\end{equation}
\end{theorem}

\begin{proof}
由留数定理，
\(
J_0=2\pi i\,\operatorname{Res}_{u=u_0}F
\)
且
\(
J_1=2\pi i\,u_0\,\operatorname{Res}_{u=u_0}F
\)，
故读出式成立。

设 \(\Delta_k:=J_k-\widetilde J_k\)。
则
\[
|\Delta_k|
\le
\oint_{|u|=R}|u|^k|F-\widetilde F|\,|\dd u|
\le
2\pi R^{k+1}\eta,
\qquad k=0,1.
\]
于是
\[
\left|\operatorname{Res}F-\frac{1}{2\pi i}\widetilde J_0\right|
=\frac{|\Delta_0|}{2\pi}\le R\eta,
\]
即 \eqref{eq:xi-contour-residue-error-bound}。
并由 \(|J_0|\ge |\widetilde J_0|-|\Delta_0|>0\) 得 \(J_0\neq 0\)。

再由分式差估计
\[
\left|\frac{J_1}{J_0}-\frac{\widetilde J_1}{\widetilde J_0}\right|
\le
\frac{|J_1-\widetilde J_1|}{|J_0|}
+
\frac{|\widetilde J_1|\,|J_0-\widetilde J_0|}{|J_0|\,|\widetilde J_0|},
\]
并代入
\[
|J_1-\widetilde J_1|\le 2\pi R^2\eta,\qquad
|J_0-\widetilde J_0|\le 2\pi R\eta,\qquad
|J_0|\ge |\widetilde J_0|-2\pi R\eta,
\]
以及
\[
|\widetilde J_1|
\le
\oint_{|u|=R}|u|\,|\widetilde F(u)|\,|\dd u|
\le
2\pi R^2\sup_{|u|=R}|\widetilde F(u)|,
\]
即可化得 \eqref{eq:xi-contour-pole-location-error-bound}。证毕。
\end{proof}

\begin{remark}[与同余块极点审计的接口]\label{rem:xi-congruence-block-contour-audit-interface}
将定理 \ref{thm:xi-single-pole-contour-readout-error-audit}
应用于 \(\mathcal A_{b,m}^{(j)}\) 在局部隔离轮廓中的单极点读出，
即可把“同余块留数锁相”转写为一条可递归缩放的误差审计链：
轮廓半径 \(R\) 与边界逼近误差 \(\eta\) 的选择直接给出
极点位置与留数的不确定度闭式上界。
\end{remark}

\endinput
